% GENERATED_BY: oc_core_monograph_translation_draft_pipeline_v1.py
% AI_ASSISTED_TRANSLATION_DRAFT: TRUE
% THEOREM_REVIEW_STATUS: PENDING
% THEOREM_REVIEW_MODE: FORMAL_AI_GATE__CODEX_CLI_CHATGPT
% THEOREM_REVIEW_REVIEWER_ID:
% THEOREM_REVIEW_AUTHORITY_CLASS:
% THEOREM_REVIEWED_AT:
% THEOREM_REVIEW_DOSSIER_REF:
% SOURCE_LANGUAGE: EN
% TARGET_LANGUAGE: RU
% SOURCE_EN_REF: content/m_spaces/m2.tex
% ================================================================
% ==== FILE: content/m_spaces/m2.tex
% ================================================================

% ==============================
%  Ontology of Continua — Core
%  Meta-Space: M2
%  FULL MODULE — FINAL
% ==============================

\subsubsection{\texorpdfstring{$M_2$}{M_2} Обзор}
\label{sec:m2-overview}

$M_2$ — это мета-пространство, задающее структуру допустимости для
двумерного континуума $K_2$.
Переход в $M_2$ соответствует возникновению:

\begin{itemize}
    \item размерности пространства за пределами $A_1$,
    \item локальной геометрической структуры,
    \item динамической метрики,
    \item причинных конусов и допустимых скоростей распространения,
    \item порогов квантизации и минимальных кластеров,
    \item полной динамики континуумов, заданной действием.
\end{itemize}

Это первое мета-пространство, где одновременно допустимы
геометрия, время и физическое распространение.

\subsubsection*{1. Структура \texorpdfstring{$\Omega(M_2)$}{\Omega(M_2)}}

Допустимое пространство конфигураций в $M_2$:

\[
\Omega(M_2)
=
\Big\{
\phi : X_2 \to V \ \big|\
\phi \in C^{0}(X_2,V),\
\partial_i\phi,\ \partial_t\phi \in H^1,\
 g_{\mu\nu} \in C^{0}
\Big\}.
\]

Требования:

\begin{enumerate}
    \item \textbf{Двумерная топология.}
    $X_2$ должен быть двумерным многообразием с локальными координатными
    картами.
    Типичные примеры: $\mathbb{R}^2$, цилиндр или компактная поверхность.

    \item \textbf{Локальная дифференцируемость.}
    Первые производные $\partial_i\phi$ и $\partial_t\phi$ должны существовать
    в соболевом смысле.

    \item \textbf{Допустимое метрическое поле.}
    $M_2$ допускает динамическую метрику $g_{\mu\nu}$, но только с сигнатурой:
    \[
    (-,+,+),
    \]
    возникающей как гессиан функционала структурной натяжки $T_2$.

    \item \textbf{Ограничение локальности.}
    Допустимые взаимодействия должны быть локальными на $X_2$.

    \item \textbf{Конечная скорость распространения.}
    Причинные конусы должны существовать (\(\Theta_{\text{conn}}\)
    обеспечивает допустимость светового конуса).
\end{enumerate}

Эти условия задают допустимую область для любого континуума $K_2$.

\subsubsection*{2. Допустимые оси в \texorpdfstring{$M_2$}{M_2}}

$M_2$ допускает ровно две независимые пространственные оси:

\[
A_1,\ \ A_2,
\]
где $A_1$ наследуется, а $A_2$ возникает из размерностного порога:
\[
T_1 \ge \Theta_{\text{dim}}.
\]

В $M_2$ не допускаются дополнительные оси.
Более высокие оси требуют переходов в $M_3$ и выше.

\subsubsection*{3. Допустимые потенциалы и энергии}

Мета-пространство $M_2$ задаёт функционалы энергии вида:

\[
E[\phi]
=
\int_{X_2}
\left(
\frac{1}{2} g^{ij} \partial_i \phi \partial_j \phi
+ V(\phi)
\right) \sqrt{|g|}\, d^2x.
\]

Допустимые потенциалы $P(M_2)$ должны удовлетворять:

\begin{itemize}
    \item коэрцитивности для устойчивости к коллапсу,
    \item ограниченности снизу,
    \item дифференцируемости для порождения потоков,
    \item исключительно локальной зависимости (недопустимы нелокальные потенциалы).
\end{itemize}

\subsubsection*{4. Пороги и структурная натяжка}

$M_2$ задаёт:

\[
\Theta_2 = \inf \{ T_2(\phi, g_{\mu\nu}) : (\phi, g) \in \Omega(M_2) \},
\]

со структурной натяжкой:

\[
T_2 = \int_{X_2}
\left(
g^{ij} \partial_i \phi \partial_j \phi
+ W(\phi)
+ \alpha R
\right)
\sqrt{|g|}\, d^2x,
\]

где:

\begin{itemize}
    \item $W(\phi)$ вводит ограничения на допустимые конфигурации,
    \item $R$ — скалярная кривизна,
    \item $\alpha$ управляет геометрической жёсткостью.
\end{itemize}

Ключевая особенность $M_2$:
\textbf{метрика $g_{\mu\nu}$ возникает из второй вариации $T_2$},
то есть геометрия является производной концепцией, а не исходной.

\subsubsection*{5. Допустимые потоки и распространение}

Допустимые потоки:

\[
J_2 = (\partial_i \phi,\ \partial_t \phi,\ \nabla_i T_2,\ \Box_g \phi),
\]

при д'Аламбертиане, определяемом через $g$.

Ограничения на распространение:

\begin{itemize}
    \item сигналы должны лежать внутри причинного конуса,
    определяемого $\Theta_{\text{conn}}$,
    \item скорость распространения не может превышать допустимый предел,
    заданный $g_{\mu\nu}$,
    \item потоки должны сохранять регулярность ($H^1$-совместимость).
\end{itemize}

Эти условия кодируют возникновение физической причинности в $K_2$.

\subsubsection*{6. Пороги квантизации и минимальные кластеры}

$M_2$ — первое мета-пространство, вводящее:

\[
\Theta_{\text{quant}} > 0.
\]

Отсюда следует:

\begin{itemize}
    \item минимальный устойчивый размер кластера $q_{\min}$,
    \item дискретные действия операторов:
    \[
    a^\dagger |k\rangle,\quad a |k\rangle,
    \]
    как изменения меры связности $k$,
    \item квантизованные возбуждения как топологически устойчивые
    конфигурации в $\Omega(M_2)$.
\end{itemize}

Таким образом, квантизация является структурной необходимостью,
а не внешней аксиомой.

\subsubsection*{7. Коллапс и граница \texorpdfstring{$M_2$}{M_2}}

Конфигурация касается $\partial\Omega(M_2)$, если:

\begin{itemize}
    \item метрика вырождена ($\det g \to 0$),
    \item нарушается локальность,
    \item $T_2 \ge \Theta_2$,
    \item кривизна расходится ($|R| \to \infty$),
    \item допустимые оси перестают быть измеримыми.
\end{itemize}

Преодоление этой границы вынуждает коллапс любого континуума $K_2$.

\subsubsection*{8. Связь с \texorpdfstring{$K_1$}{K_1} и \texorpdfstring{$K_3$}{K_3}}

$M_2$ делает переход $K_1 \to K_2$ возможным через $\Psi_{1\to2}$:

\[
A_1 \mapsto (A_1, A_2),\quad
\Omega(K_1) \to \Omega(K_2),\quad
T_1 \to T_2,\quad
\text{dim} = 2.
\]

Переход в $K_3$ допустим только при:

\[
A_3 \in A(M_3),\qquad
T_2 \ge \Theta_{\text{dim}},
\qquad
\Omega(K_3) \subseteq \Omega(M_3).
\]

Следовательно, $M_2$ — это мета-пространство:

\begin{itemize}
    \item геометрии,
    \item причинности,
    \item квантизации,
    \item размерностной устойчивости,
    \item распространения.
\end{itemize}

Оно составляет структурную основу всех более высоких континуумов.
